\scanpage{82}
\begin{center}
INSTITUT\\
DES HAUTES ÉTUDES SCIENTIFIQUES\\[-.3ex]
\rule{1.3cm}{.4pt}
\end{center}

Appendice à $\underline{\overline{\mathrm{XXI}}}$ de Katz : torsion en coh. cristalline

$\underline{\overline{\mathrm{XXII}}}$ \quad Référence à la démonstration de Deligne du th. de Congruences de Katz (\unsure{généralisée|génér. lisse})

{\reversemarginpar\margin{$\underline{\overline{\unsure{\mathrm{VI}|\mathrm{XI}}}}$\\ou $\underline{\overline{\mathrm{XIII}}}$\unsure{?|'}}}Vérifier qu'on y trouve le fait qu'en déformant une singularité quadratique ordinaire, \unsure{---|on} \unsure{donne} au plus autant (y inclus cas de car. rés. 2 et dim relative paire). En d'autres termes, l'ens des pts où un morphisme plat \deleted{de}\inserted{loc.} prés. finie est lisse ou a une singularité quadratique ordinaire est tjrs. \uline{ouvert}.
% archive sheet number: 81
